%! arara: clean: {
%! arara: --> extensions:
%! arara: --> ['aux', 'bbl', 'bcf', 'blg', 'log', 'nav', 'out', 'pdf', 'run.xml', 'snm', 'toc']
%! arara: --> }
%! arara: lualatex: {
%! arara: --> shell: yes,
%! arara: --> draft: yes,
%! arara: --> interaction: batchmode
%! arara: --> }
%! arara: biber
% arara: lualatex: {
% arara: --> shell: yes,
% arara: --> draft: no,
% arara: --> interaction: batchmode
% arara: --> }
% arara: lualatex: {
% arara: --> shell: yes,
% arara: --> draft: no,
% arara: --> interaction: batchmode
% arara: --> }
%! arara: clean: {
%! arara: --> extensions:
%! arara: --> ['aux', 'bbl', 'bcf', 'blg', 'log', 'nav', 'out', 'run.xml', 'snm', 'toc']
%! arara: --> }
\input{preamble}

\pointpoints{Punto}{Puntos}
\bracketedpoints
\begin{document}
\begin{center}
	\sffamily\bfseries\Large
	Métodos Numéricos para Ecuaciones Diferenciales Parciales~CM 032 A\\
	Laboratorio N$^{\circ}1$
\end{center}

\begin{center}
	\fbox{\fbox{\parbox{5.5in}
			{\centering
				Responda las preguntas en los espacios provistos en las hojas
				de preguntas.
				Los argumentos y la claridad de las respuestas se
				considerarán en la puntuación final.
			}}}
\end{center}
\vspace{0.1in}
\makebox[\textwidth]{Nombres y apellidos:\enspace\hrulefill}
\vspace{0.2in}
\makebox[\textwidth]{Nombres y apellidos del profesor:\enspace Fidel Jara Huanca.\hfill}

Estudie el material \textcolor{red}{\href{https://sbu-python-class.github.io/python-science}{Python for Scientific Computing}}
y responda las siguientes preguntas.

\begin{questions}
	\question

	Cuando hablamos de punto flotante, discutimos el
	\textit{épsilon de la máquina}, $\epsilon$ -este es el número más
	pequeño que cuando se suma a $1$ sigue siendo diferente de $1$.
	Calcularemos $\epsilon$ aquí:
	\begin{itemize}
		\item

		      Elija una estimación inicial para $\epsilon$ de
		      \mintinline{python}|eps = 1|.

		\item

		      Cree un bucle que compruebe si
		      \mintinline{python}|1 + eps| es diferente de
		      \mintinline{python}|1|.

		\item

		      En cada iteración del bucle, corte el valor de
		      \mintinline{python}|eps| por la mitad.
	\end{itemize}
	¿Qué valor de $\epsilon$ encuentras?

	\question

	Para iterar sobre las tuplas, donde la \textit{i}ésima tupla contiene
	los \textit{i}-ésimos elementos de ciertas secuencias, podemos usar la función
	\mintinline{python}|zip(*secuencias)|.

	Iteraremos sobre dos listas, \mintinline{python}|nombres| y
	\mintinline{python}|edad| e imprima las tuplas resultantes.

	\begin{itemize}
		\item

		      Empiece por inicializar las listas
		      \mintinline{python}|nombres = ["Mary", "John", "Sarah"]| y
		      \mintinline{python}|edad = [21, 56, 98]|.

		\item


		      Iterar sobre las tuplas que contienen un nombre y una edad,
		      la función \mintinline{python}|zip(list1, list2)| la
		      función podría ser útil aquí.

		\item

		      Imprima cadenas formateadas del tipo
		      \mintinline{text}|"Nombre tiene EDAD años"|.
	\end{itemize}

	\question

	La función \mintinline{python}|enumerate(secuencia)| retorna
	tuplas que contienen índices de objetos en la secuencia, y los
	objetos.

	El módulo \mintinline{python}|random| proporciona herramientas para
	trabajar con los números aleatorios.
	En particular, \mintinline{python}|random.randint(start, end)|
	genera un número aleatorio no menor que
	\mintinline{python}|start|, y no mayor que \mintinline{python}|end|.

	\begin{itemize}
		\item

		      Genere una lista de $10$ números aleatorios de $0$ a $9$.

		\item

		      Usando la función \mintinline{python}|enumerate(random_list)|,
		      iterar sobre las tuplas de números aleatorios y sus índices,
		      e imprimirlos \mintinline{text}|"Coincidencia: NÚMERO e ÍNDICE"|
		      si el número aleatorio y su índice en la lista coinciden.
	\end{itemize}
	% \inputminted[firstline=1,lastline=43,fontsize=\tiny,linenos,highlightlines={23-28}]{scilab}{main.sci}
	\inputminted[fontsize=\tiny,linenos]{python}{q3.py}

	\question

	La secuencia de Fibbonacci es una secuencia numérica donde cada
	número es la suma de los números anteriores $2$, por ejemplo,
	$1$, $1$, $2$, $3$, $5$, $8$, $13$, \lpuntos.
	Crea una lista donde los elementos son los términos de la secuencia
	de Fibbonacci:
	\begin{itemize}
		\item

		      Comience con la lista \mintinline{python}|fib = [1, 1]|.

		\item

		      Repita $25$ veces, calcule el siguiente término como la suma de
		      términos de $2$ anteriores y añádalos a la lista.

		\item

		      Una vez completado el ciclo, imprima los términos.
	\end{itemize}
	Puede resultarle útil utilizar \mintinline{python}|fib[-1]| y
	\mintinline{python}|fib[-2]| para acceder a los dos últimos elementos del
	lista.

	\question

	Podemos usar la función \mintinline{python}|input()| para preguntar por
	la entrada desde el prompt.

	Cree una lista vacía y use un bucle \mintinline{python}|while| para solicitar
	información al usuario y agregue sus entradas a la lista.
	Continúe en bucle hasta que se agreguen $10$ elementos a la lista.

	\question

	Here is a list of book titles
	(from \url{http://thegreatestbooks.org}).
	Loop through the list and capitalize each word in each title.
	You might find the \mintinline{python}|.capitalize()| method that
	works on strings useful.
	\inputminted[fontsize=\tiny,linenos]{python}{q6.py}

	\question

	Here's some text (the Gettysburg Address).
	Our goal is to count how many times each word repeats.
	We'll do a brute force method first, and then we'll look a ways to
	do it more efficiently (and compactly).

	\inputminted[fontsize=\tiny,linenos]{python}{q7.py}

	We've already seen the \mintinline{python}|.split()| method will,
	by default, split by spaces, so it will split this into words,
	producing a list:

	Now, the next problem is that some of these still have punctuation.
	In particular, we see \mintinline{python}|"."|,
	\mintinline{python}|","|, and \mintinline{python}|"--"|.

	When considering a word, we can get rid of these by using the
	\mintinline{python}|.replace()| method:

	Another problem is case-we want to count ``but'' and ``But'' as the same.
	Strings have a \mintinline{python}|lower()| method that can be used
	to convert a string:
	Recall that strings are immutable, so \mintinline{python}|`replace()`| produces a new string on output.

	Create a dictionary that uses the unique words as keys and has as a
	value the number of times that word appears.

	Write a loop over the words in the string (using our split version)
	and do the following:
	\begin{itemize}
		\item

		      Remove any punctuation.

		\item

		      Convert to lowercase.

		\item

		      Test if the word is already a key in the dictionary (using the `in` operator).

		\item

		      If the key exists, increment the word count for that key.

		\item

		      Otherwise, add it to the dictionary with the appropriate count of \mintinline{python}|1|.
	\end{itemize}

	At the end, print out the words and a count of how many times they appear.

	\question

	We can actually do this a lot more compactly by using another list
	comprehensions and another python datatype called a set.
	A set is a group of items, where each item is unique (e.g., no
	repetitions).
	Here's a list comprehension that removes all the punctuation and converts to lower case:
	and by using the \mintinline{python}|set()| function, we turn the list into a set, removing any duplicates:
	now we can loop over the unique words and use the \mintinline{python}|count()| method of a list to find how many there are
	Even shorter -- we can use a dictionary comprehension, like a list comprehension.

	\inputminted[fontsize=\tiny,linenos]{python}{q8.py}

	\question

	Let's practice functions.
	Here's a simple function that takes a string and returns a list of
	all the $4$ letter words:

	Write a version of this function that takes a second argument, $n$,
	that is the word length we want to search for.

	\question

	A prime number is divisible only by 1 and itself.
	We want to write a function that takes a positive integer, $n$,
	and finds all of the primes up to that number.

	A simple (although not very fast) way to find the primes is to
	start at $1$, and build a list of primes by checking if the current
	number is divisible by any of the previously found primes.
	If it is not divisible by any earlier primes, then it is a prime.

	The modulus operator, \% could be helpful here.

	\question

	We want to safely convert a string into a \mintinline{python}|float|,
	\mintinline{python}|int|, or leave it as a string, depending on its contents.
	As we've already seen, python provides \mintinline{python}|float()|
	and \mintinline{python}|int()| functions for this:

	Notice that an \mintinline{python}|int| can be converted to a
	\mintinline{python}|float|, but if you convert a
	\mintinline{python}|float| to an \mintinline{python}|int|, you rise
	losing significant digits.
	A string cannot be converted to either.

	Write a function, \mintinline{python}|convert_type(a)| that takes a
	string \mintinline{python}|a|, and converts it to a float if it is
	a number with a decimal point, an int if it is an integer, or
	leaves it as a string otherwise, and returns the result.
	You'll want to use exceptions to prevent the code from aborting.

	\inputminted[fontsize=\tiny,linenos]{python}{q10.py}

	\question

	Here we'll write a simple tic-tac-toe game that 2 players can play.
	First we'll create a string that represents our game board:

	This board will look a little funny if we just print it—the spacing is set to look right when we replace the {} with x or o

	and well use a dictionary to denote the status of each square, “x”, “o”, or empty, “”

	Note that our {} placeholders in the board string have identifiers (the numbers in the {}).
	We can use these to match the variables we want to print to the placeholder in the string,
	regardless of the order in the format()
	Here’s an easy way to add the values of our dictionary to the appropriate squares in our game board. First note that each of the {} is labeled with a number that matches the keys in our dictionary.
	Python provides a way to unpack a dictionary into labeled arguments, using **
	This lets us to write a function to show the tic-tac-toe board.

	Now we need a function that asks a player for a move:

	Using the functions defined above,

	\begin{itemize}
		\item \mintinline{python}|initialize_board()|
		\item \mintinline{python}|show_board()|
		\item \mintinline{python}|get_move()|
	\end{itemize}
	fill in the function \mintinline{python}|play_game()| below to complete the game,
	asking for the moves one at a time, alternating between player $1$ and $2$.

	\question

	Let's write a simple shopping cart class - this will hold items
	that you intend to purchase as well as the amount, etc.
	And allow you to add / remove items, get a subtotal, etc.

	We'll use two classes:
	Item will be a single item and \mintinline{python}|ShoppingCart|
	will be the collection of items you wish to purchase.

	First, our store needs an inventory - here's what we have for sale:

	Here's the start of an item class - we want it to hold the name and quantity.

	You should have the following features:

	\begin{itemize}
		\item

		      the name should be something in our inventory.

		\item

		      Our shopping cart will include a list of all the items we
		      want to buy, so we want to be able to check for duplicates.
		      Implement the equal test,\mintinline{python}|==|, using
		      \mintinline{python}|__eq__|.

		\item

		      We'll want to consolidate dupes, so implement the
		      \mintinline{python}|+| operator, using
		      \mintinline{python}|__add__| so we can add items together in our shopping cart.
		      Note, add should raise a \mintinline{python}|ValueError| if you try to add two \mintinline{python}|Items| that don't have the same name.
	\end{itemize}
	Here's a start:
\end{questions}
\vfill
\begin{flushright}\bfseries
	Facultad de Ciencias\\[2mm]
	Universidad Nacional de Ingeniería\\[2mm]
	25 de marzo de 2026
\end{flushright}
\end{document}
